\chapter{The Sitas Executor}

The sitas executor is the userspace runtime that drives futures on one shard. This
chapter explains how it works, how it interacts with the completion queue, and
why the unified wait is the central insight of the co-design.

\section{One executor per shard}

Each sitas shard has:

\begin{itemize}
    \item \textbf{One executor} --- a cooperative scheduler that polls ready
        futures within a budget.
    \item \textbf{One ready queue} --- tasks whose futures returned
        \texttt{Poll::Pending} and whose wakers have since been invoked.
    \item \textbf{One timer wheel} --- futures awaiting deadlines.
    \item \textbf{One CQ partition} --- the completion queue on which this shard
        waits (\texttt{CQ\_WAIT}).
    \item \textbf{Shard-owned service state} --- mutable data that only this
        executor touches.
\end{itemize}

The executor runs on one LP-affine thread. When that thread is scheduled onto its
LP by the kernel, the executor runs. When the thread blocks (no ready tasks), the
LP is free to run another domain's thread.

\section{The polling loop}

The executor's main loop is:

\begin{verbatim}
loop {
    // 1. Poll ready tasks within a budget.
    while budget_remaining > 0 && ready_tasks.not_empty() {
        let task = ready_tasks.pop();
        match task.future.poll(task.waker) {
            Poll::Ready(output) => { /* task done */ }
            Poll::Pending => { /* will be re-awoken by waker */ }
        }
        budget_remaining -= 1;
    }

    // 2. If tasks remain but budget exhausted, re-queue and yield.
    if ready_tasks.not_empty() {
        ready_tasks.requeue_for_next_tick();
    }

    // 3. No ready tasks? Block in the unified wait.
    if ready_tasks.is_empty() && timer_wheel.is_empty() {
        reactor.wait_for_events(); // -> CQ_WAIT
    }
}
\end{verbatim}

The budget prevents one task (or a tight loop of always-ready tasks) from
starving other tasks on the shard. If the budget is exhausted with work remaining,
the executor re-queues the remaining tasks and yields (either to other threads on
the same LP via the kernel scheduler, or to idle).

\section{The unified wait --- three event sources, one wait point}

When the executor has no ready tasks and no pending timers, it calls
\texttt{CQ\_WAIT} on the shard's completion queue. This is the \textbf{only}
blocking point in the shard's execution.

Three event sources converge on this one wait point:

\subsection{Kernel completions (CQ entries)}

When a submitted operation completes (timer expired, DMA finished, device
interrupt processed), the kernel posts a \texttt{CompletionEntry} to the shard's
CQ ring and wakes the blocked thread. The executor drains the ring:

\begin{verbatim}
fn drain_cq(&mut self) {
    while let Some(entry) = self.cq_ring.read() {
        let waker = self.waker_registry.take(entry.user_data);
        if let Some(waker) = waker {
            waker.wake(); // task becomes ready
        }
    }
}
\end{verbatim}

\subsection{Endpoint readiness (coalesced wake)}

An endpoint can be bound to a CQ. When a message arrives on the endpoint, the
kernel posts a coalesced wake to the bound CQ. The executor, on waking, calls the
endpoint's receive path to drain the pending messages.

Coalescing means: if 10 messages arrive before the shard wakes, the kernel
delivers one wake, not 10. The executor drains all 10 messages in one batch.

\subsection{Device interrupts (coalesced wake)}

A device interrupt object can be bound to a CQ. When the interrupt fires, the
kernel posts a coalesced wake. The executor, on waking, calls the device's
interrupt handler (or the driver's polling function).

\section{Task wakeup --- from CQ entry to future poll}

When the executor drains a CQ entry, it looks up the registered waker for that
entry's \texttt{user\_data} cookie. The cookie was set by the task when it
submitted the operation. The waker marks the task ready:

\begin{verbatim}
// Inside a task's Future::poll:
fn poll(mut self: Pin<&mut Self>, cx: &mut Context<'_>) -> Poll<u64> {
    match self.cap.poll() {
        Ok(Some(completed)) => Poll::Ready(completed.result),
        Ok(None) => {
            // Register waker so we are notified when this completes.
            self.cap.register_waker(cx.waker().clone());
            Poll::Pending
        }
        Err(e) => Poll::Ready(/* error */),
    }
}
\end{verbatim}

The waker is a \texttt{ShardWaker}: when invoked, it sets the task's ready flag
and signals the reactor (if the reactor is currently blocked) to re-enter the
polling loop. The reactor then polls the task, which finds its completion ready.

\section{Cross-shard messaging without the kernel}

Within one protection domain, tasks on different shards communicate through sitas
typed mailboxes --- not through kernel endpoint IPC:

\begin{verbatim}
// Shard A: send work to Shard B.
let sender: ShardSender<Command> = shard_b.mailbox();
sender.try_send(Command::Process { data: owned_buffer })?;

// Shard B: receive work.
let receiver: ShardReceiver<Command> = shard_b.receiver();
loop {
    let cmd = receiver.recv().await?;
    match cmd {
        Command::Process { data } => { /* handle */ }
    }
}
\end{verbatim}

The \texttt{ShardSender/ShardReceiver} pair is a bounded, typed, MPSC channel
within userspace. Sending does not cross the kernel boundary. The receiving task
registers its waker on the channel; when a message arrives, the waker fires and
the task re-enters the executor's polling loop.

This is \textbf{internal} cross-shard communication (same protection domain,
different shards). External communication (different protection domains) uses
endpoint IPC (Chapter~5).

\section{The ShardParker --- blocking on native state}

When the executor has no ready tasks, it must block its OS thread. The mechanism
is the \texttt{ShardParker} trait:

\begin{verbatim}
trait ShardParker {
    fn park(&self);    // Block the current thread.
    fn unpark(&self);  // Wake the blocked thread (from another shard).
}
\end{verbatim}

On CharlotteOS, the implementation is:

\begin{verbatim}
fn park(&self) { syscall::cq_wait(self.cq_id); }
fn unpark(&self) { syscall::cq_wake(self.cq_id, self.target_lp); }
\end{verbatim}

The parker maps directly to \texttt{CQ\_WAIT} / \texttt{CQ\_WAKE}. There is no
intermediate thread pool, no \texttt{pthread\_cond\_wait}, and no busy-loop.
The thread parks on native completion and message state.

This is the heart of the ``no retrofit tax'' claim: the sitas runtime does not
emulate async on top of a blocking kernel because the kernel's blocking primitive
(\texttt{CQ\_WAIT}) \emph{already} understands the semantics of ``wait for any
of several async sources.'' The executor doesn't need to build a reactor on top of
\texttt{epoll} or a thread pool; the kernel \emph{is} the reactor.

\section{Why the sitas executor matters for critical software}

\begin{enumerate}
    \item \textbf{Cooperative scheduling means no data races on shard-local state.}
        Within a shard, only one task runs at a time. The executor polls tasks
        sequentially. There is no preemption within the shard. This means
        shard-local mutable state never needs locks.

    \item \textbf{The single wait point eliminates a class of race conditions.}
        The executor does not wait on ``timer or CQ or endpoint'' through three
        separate blocking calls with three separate check-and-arm sequences. One
        CQ, one wait, one drain loop. The kernel guarantees the wake is not lost
        between checking and arming.

    \item \textbf{No thread-pool overhead.} The executor uses exactly one kernel
        thread per shard. When that thread blocks, it blocks on the CQ --- the
        same object that will deliver its wake. There is no need for a pool of
        threads to convert blocking syscalls into futures.

    \item \textbf{Budgeting prevents starvation.} The polling budget ensures that
        a shard with many ready tasks makes progress on all of them, not just the
        one that finishes fastest. This prevents priority-inversion-like behavior
        within a shard.

    \item \textbf{Internal messaging avoids unnecessary kernel crossings.} Shards
        within the same process communicate through typed channels in userspace.
        Only cross-process IPC crosses the kernel boundary. This keeps the fast
        path fast and the trusted computing base small.
\end{enumerate}

\endinput
